% =================================================================
% ==================== CONTINUATION: Q251--Q350 ===================
% =================================================================
% =================================================================
% ==================== CONTINUATION: Q151--Q250 ===================
% =================================================================
\documentclass[aspectratio=169,8pt]{beamer}
\usetheme{Madrid}
\usepackage[utf8]{inputenc}
\usepackage[T1]{fontenc}
\usepackage{amsmath,amssymb}
\usepackage{booktabs}
\usepackage{enumitem}
\setlist[enumerate]{label=\Alph*.,leftmargin=*,itemsep=1pt,topsep=2pt}
\setlist[itemize]{leftmargin=*,itemsep=1pt,topsep=2pt}
\setbeamerfont{block body}{size=\tiny}
\setbeamerfont{block title}{size=\scriptsize}
\setbeamerfont{frametitle}{size=\footnotesize}
\setbeamerfont{normal text}{size=\tiny}
\setbeamertemplate{navigation symbols}{}
\setbeamertemplate{itemize item}{\tiny$\bullet$}
\setbeamertemplate{itemize subitem}{\tiny$\circ$}

\title[GARP RAI]{GARP Risk \& AI --- 250 Detailed Hard Questions}
\subtitle{Unofficial Exam Prep}
\author{Unofficial}
\date{\today}

\begin{document}
\section{Extended Practice: Modules 1--5 (Hard)}

% ---------------- MODULE 1 — AI AND RISK ----------------

\begin{frame}{Q251 --- AI Risk Taxonomy Applied}
\begin{block}{Question}
A global asset manager deploys an LLM-based portfolio commentary generator. After release, journalists discover that the model occasionally invents fictitious analyst opinions attributed to real people at the firm. The firm's legal team, PR team, and regulators all become involved. Which combination of risk categories from the GARP taxonomy does this scenario most directly implicate?
\end{block}
\begin{enumerate}
\item Individual risk and operational risk only.
\item Organizational risk (reputational and legal) and societal risk (misinformation at scale).
\item Societal risk only.
\item Strategic risk only.
\end{enumerate}
\begin{exampleblock}{Answer: B}\end{exampleblock}
\begin{block}{Explanation}
Reputational damage, legal liability, and regulatory scrutiny are organizational risks. The generation of misleading synthetic content attributed to real people amplifies misinformation risk at the societal level. Option A understates the scope. Option C misses the firm-level harm. Option D is too narrow.
\end{block}
\end{frame}

\begin{frame}{Q252 --- Human Oversight in Critical Decisions}
\begin{block}{Question}
A bank's automated credit-approval system generates a recommended decision for every application. The bank has decided that for applications with a predicted default probability between 15\% and 25\% (an uncertain band), the model's recommendation should be reviewed by a human underwriter. What is the most appropriate rationale for this design choice?
\end{block}
\begin{enumerate}
\item To reduce compute costs.
\item To ensure meaningful human oversight in a zone of genuine uncertainty, catching model errors and exercising judgment where the model is not confident.
\item Because the regulator requires it for all applications.
\item Because humans are always more accurate than models.
\end{enumerate}
\begin{exampleblock}{Answer: B}\end{exampleblock}
\begin{block}{Explanation}
Human-in-the-loop review is best used in zones of genuine model uncertainty, where the marginal value of human judgment is greatest. Option A is a secondary benefit. Option C is over-broad. Option D is an overstatement.
\end{block}
\end{frame}

\begin{frame}{Q253 --- Deep Learning vs Traditional Econometrics}
\begin{block}{Question}
A quant team needs to build a daily volume forecast for a large-cap equity. Two approaches are proposed: (i) a traditional ARIMA model, and (ii) a recurrent neural network (LSTM) trained on high-frequency order-book data. The team has 8 years of minute-level data. Which approach is most likely to be appropriate, and why?
\end{block}
\begin{enumerate}
\item ARIMA, because neural networks cannot handle sequential data.
\item LSTM, because it can capture complex non-linear dependencies in high-frequency data that linear time-series models may miss, provided sufficient data and computational resources.
\item ARIMA, because it is always more accurate for financial time series.
\item Neither approach; only tree-based models apply.
\end{enumerate}
\begin{exampleblock}{Answer: B}\end{exampleblock}
\begin{block}{Explanation}
LSTMs are specifically designed for sequential data and can capture non-linear patterns in high-frequency series. ARIMA remains useful but is linear and limited. Option A is wrong. Option C is not generally true. Option D is incorrect.
\end{block}
\end{frame}

\begin{frame}{Q254 --- Data Visualization for Risk Detection}
\begin{block}{Question}
A risk analyst generates a scatter plot of trading volumes versus trade sizes for a broker's equity order flow. A small cluster of orders shows unusually high volume with unusually small trade size, far from the main cloud of observations. What might this cluster indicate?
\end{block}
\begin{enumerate}
\item Normal market activity.
\item A potential anomaly or pattern of concern, such as algorithmic order splitting or spoofing behaviour, warranting further investigation.
\item A data-entry error with certainty.
\item Nothing of interest.
\end{enumerate}
\begin{exampleblock}{Answer: B}\end{exampleblock}
\begin{block}{Explanation}
Visual inspection of features can reveal clusters or outliers that indicate anomalies worth investigating. Option A ignores the deviation. Option C is premature. Option D misses the signal.
\end{block}
\end{frame}

\begin{frame}{Q255 --- Data Cleaning Trade-offs}
\begin{block}{Question}
A junior analyst decides to remove all observations with any missing value from a dataset used for a credit-scoring model. The dataset originally had 100,000 rows; after the removal, only 62,000 remain. The senior data scientist raises a concern. What is the most likely concern?
\end{block}
\begin{enumerate}
\item The dataset is too small.
\item Listwise deletion can introduce bias if missingness is systematic, distorting the model's representation of the underlying population.
\item The dataset should be expanded.
\item The dataset cannot be used.
\end{enumerate}
\begin{exampleblock}{Answer: B}\end{exampleblock}
\begin{block}{Explanation}
Listwise deletion is problematic when missingness is non-random. Biases can be introduced if missingness correlates with the target or other important features. Option A ignores the deeper issue. Options C and D are secondary or incorrect.
\end{block}
\end{frame}

\begin{frame}{Q256 --- Feature Engineering in Credit Models}
\begin{block}{Question}
A credit risk team considers adding a new feature to a default model: ``number of missed payments in the last 6 months.'' The team debates whether this feature should be used at training and prediction time. Which is the best approach?
\end{block}
\begin{enumerate}
\item Use it at training only.
\item Use it at both training and prediction if it is available at prediction time in real-world deployment; otherwise it introduces leakage.
\item Use it at prediction only.
\item Do not use it at all.
\end{enumerate}
\begin{exampleblock}{Answer: B}\end{exampleblock}
\begin{block}{Explanation}
Features must be available at both training and prediction time. Features computed using future information cause leakage. Options A, C, and D are incorrect.
\end{block}
\end{frame}

\begin{frame}{Q257 --- Regularization in Practice}
\begin{block}{Question}
A team fits a logistic regression with L2 penalty to a credit-default dataset with 5,000 observations and 200 features. They observe that increasing the regularization strength decreases both training and test accuracy. What does this suggest?
\end{block}
\begin{enumerate}
\item The regularization strength should be further increased.
\item The model may be underfitting; regularization strength is too high, or the model needs more complexity/features.
\item Regularization is not needed.
\item The target variable is wrong.
\end{enumerate}
\begin{exampleblock}{Answer: B}\end{exampleblock}
\begin{block}{Explanation}
If increasing regularization decreases both training and test performance, the model is being over-penalised (bias dominates). The right regularization strength balances bias and variance. Options A, C, and D are not supported by the symptom.
\end{block}
\end{frame}

\begin{frame}{Q258 --- PCA and Interpretability Trade-off}
\begin{block}{Question}
An analyst applies PCA to 20 correlated macro features and retains 4 components capturing 92\% of variance. The resulting model performs well but the business team complains that the components are hard to interpret. Which trade-off is at play?
\end{block}
\begin{enumerate}
\item Accuracy vs cost.
\item Dimensionality reduction and stability vs interpretability --- components are linear combinations of features, obscuring business meaning.
\item Overfitting vs underfitting.
\item Randomness vs determinism.
\end{enumerate}
\begin{exampleblock}{Answer: B}\end{exampleblock}
\begin{block}{Explanation}
PCA improves dimensionality and stability but sacrifices interpretability, since components do not correspond to individual business features. Options A, C, and D are not the primary trade-off.
\end{block}
\end{frame}

\begin{frame}{Q259 --- Tree-Based Feature Importance}
\begin{block}{Question}
A random forest model ranks ``age'' as the most important feature for predicting fraud. A careful analyst notices that ``age'' and ``retirementstatus'' are highly correlated. What is a known limitation of feature importance in this context?
\begin{enumerate}
\item Feature importance is always wrong.
\item Feature importance can be inflated for correlated features; the model may arbitrarily select one and mask the importance of the correlated sibling.
\item Feature importance is only valid for linear models.
\item Feature importance is a fairness metric.
\end{enumerate}
\end{block}
\begin{exampleblock}{Answer: B}\end{exampleblock}
\begin{block}{Explanation}
Impurity-based feature importance is biased toward correlated features. Permutation importance or SHAP is often more reliable in the presence of correlation. Options A, C, and D are incorrect.
\end{block}
\end{frame}

\begin{frame}{Q260 --- Data Transformation and Interpretability}
\begin{block}{Question}
A model uses log-transformed revenue as a feature. The coefficient on log(revenue) is 0.30. How should this coefficient be interpreted?
\end{block}
\begin{enumerate}
\item A one-unit increase in revenue increases the target by 0.30.
\item A 1\% increase in revenue increases the target by approximately 0.30\% (if target is also in logs) or by 0.003 units (if target is in levels), depending on the model specification.
\item The coefficient is meaningless.
\item Log transformation destroys interpretability.
\end{enumerate}
\begin{exampleblock}{Answer: B}\end{exampleblock}
\begin{block}{Explanation}
In log-log models, coefficients are elasticities. In log-linear models, coefficients are semi-elasticities. Option A misses the log scale. Options C and D are incorrect.
\end{block}
\end{frame}

\begin{frame}{Q261 --- The Bias-Variance Trade-off in Neural Nets}
\begin{block}{Question}
A deep neural network with 10 million parameters achieves 99.5\% training accuracy and 68\% test accuracy on a labelled financial dataset of 100,000 observations. What should the team do?
\end{block}
\begin{enumerate}
\item Increase model size.
\item Reduce model complexity, add regularization (dropout, weight decay), gather more data, or use early stopping --- the model is overfitting.
\item Use a simpler loss function.
\item Change the target variable.
\end{enumerate}
\begin{exampleblock}{Answer: B}\end{exampleblock}
\begin{block}{Explanation}
The huge gap between training and test accuracy signals overfitting. Standard remedies: reduce complexity, regularize, gather more data. Options A, C, and D do not address overfitting.
\end{block}
\end{frame}

\begin{frame}{Q262 --- Interpretability vs Accuracy in Regulated Environments}
\begin{block}{Question}
In a regulated credit-scoring context, a bank must explain individual decisions to customers and regulators. The bank's data science team proposes a highly accurate gradient-boosted model. Which combination of strategies best satisfies both accuracy and explainability needs?
\end{block}
\begin{enumerate}
\item Use the model as-is without explanation.
\item Use the model with post-hoc explainability (SHAP, counterfactuals) plus a simpler interpretable benchmark model for comparison and documentation.
\item Use only a logistic regression regardless of accuracy loss.
\item Use a black-box model and refuse to explain.
\end{enumerate}
\begin{exampleblock}{Answer: B}\end{exampleblock}
\begin{block}{Explanation}
Regulated environments commonly combine high-performance models with post-hoc explainability and interpretable benchmarks. Options A and D are non-compliant. Option C sacrifices accuracy without an explanation framework.
\end{block}
\end{frame}

\begin{frame}{Q263 --- Data Scaling Pitfalls}
\begin{block}{Question}
A data scientist standardizes features using the full dataset, including test data, and then splits into train/test. Another analyst says this is wrong. Which is the correct critique?
\end{block}
\begin{enumerate}
\item It is fine because standardization is a linear transformation.
\item It leaks test-set information (mean and standard deviation) into the training process, biasing performance estimates.
\item Standardization should only be applied to categorical features.
\item Standardization is unnecessary.
\end{enumerate}
\begin{exampleblock}{Answer: B}\end{exampleblock}
\begin{block}{Explanation}
Scaling parameters must be estimated on the training set only. Using the full dataset leaks test information. Options A, C, and D misstate the issue.
\end{block}
\end{frame}

\begin{frame}{Q264 --- Non-Linearity in Financial Models}
\begin{block}{Question}
A risk analyst wants to capture the relationship between interest rates and bond returns, which is known to be non-linear (convexity). Which modeling approach is most appropriate?
\end{block}
\begin{enumerate}
\item A linear regression of returns on rates.
\item A model that explicitly allows for non-linearity, such as a polynomial regression, spline, tree-based model, or neural network.
\item A logistic regression.
\item A KNN classifier.
\end{enumerate}
\begin{exampleblock}{Answer: B}\end{exampleblock}
\begin{block}{Explanation}
Non-linear relationships require models that can capture curvature: polynomial regressions, splines, trees, or neural networks. Options A, C, and D are inappropriate.
\end{block}
\end{frame}

\begin{frame}{Q265 --- Overfitting in Small Datasets}
\begin{block}{Question}
A research team fits a deep neural network to only 500 observations with 100 features. Test performance is poor. What is the most likely explanation, and what should be done?
\end{block}
\begin{enumerate}
\item The model is too simple; increase depth.
\item Overfitting due to small sample and high dimensionality; use simpler models, regularization, dimensionality reduction, or gather more data.
\item The test set is corrupted.
\item The target is mislabelled.
\end{enumerate}
\begin{exampleblock}{Answer: B}\end{exampleblock}
\begin{block}{Explanation}
With 500 observations and 100 features, a deep network will overfit. Remedies: simpler models, regularization, PCA, more data. Options A, C, and D are unlikely explanations.
\end{block}
\end{frame}

\begin{frame}{Q266 --- Missing Data in Time Series}
\begin{block}{Question}
A time-series dataset of daily FX rates has occasional gaps (weekends, holidays, occasional missing days). Which imputation approach is generally most appropriate?
\begin{enumerate}
\item Drop all rows with missing values.
\item Forward-fill, interpolation, or model-based imputation that respects temporal structure.
\item Replace with zeros.
\item Replace with the global mean.
\end{enumerate} \end{block}
\begin{exampleblock}{Answer: B}\end{exampleblock}
\begin{block}{Explanation}
Time-series missingness should be handled with methods that respect temporal order (forward-fill, interpolation, state-space models). Options A, C, and D distort the series.
\end{block}
\end{frame}

\begin{frame}{Q267 --- Categorical Encoding Pitfalls}
\begin{block}{Question}
A dataset has a high-cardinality categorical variable ``merchantid'' with 10,000 unique values. One-hot encoding would produce 10,000 new columns. What is a better approach?
\begin{enumerate}
\item Use one-hot encoding anyway.
\item Use target encoding (with proper cross-validation to avoid leakage), hashing, or embedding-based representations.
\item Remove the feature.
\item Convert to ordinal codes.
\end{enumerate} \end{block}
\begin{exampleblock}{Answer: B}\end{exampleblock}
\begin{block}{Explanation}
High-cardinality categoricals should use target encoding (with regularization), hashing, or embeddings. Option A causes dimension explosion; C loses information; D implies an order that does not exist.
\end{block}
\end{frame}

\begin{frame}{Q268 --- Outliers and Model Choice}
\begin{block}{Question}
A portfolio returns dataset has a few extreme observations due to flash-crash-type events. Two models are being compared: OLS regression and quantile regression. Which is generally more robust to these extremes?
\begin{enumerate} 
\item OLS, because it uses all information.
\item Quantile regression, because it targets conditional quantiles rather than the mean and is therefore less influenced by extreme observations.
\item Both are equally robust.
\item Neither; only tree models work.
\end{enumerate} \end{block}
\begin{exampleblock}{Answer: B}\end{exampleblock}
\begin{block}{Explanation}
Quantile regression estimates conditional quantiles, providing robustness to extreme values. OLS is sensitive. Options A, C, and D are incorrect.
\end{block}
\end{frame}

\begin{frame}{Q269 --- Testing ML Pipelines}
\begin{block}{Question}
A team deploys a production ML pipeline. Which of the following testing categories is essential to ensure reliability?
\begin{enumerate}
\item Unit tests on transformation functions.
\item Integration tests on the full pipeline.
\item Adversarial and stress tests on model behaviour.
\item All of the above.
\end{enumerate} \end{block}
\begin{exampleblock}{Answer: D}\end{exampleblock}
\begin{block}{Explanation}
Production ML pipelines need unit, integration, and adversarial/stress testing. All three are essential.
\end{block}
\end{frame}

\begin{frame}{Q270 --- Feature Selection in Practice}
\begin{block}{Question}
A team reduces 300 candidate features to 25 using LASSO, then fits a logistic regression on the selected 25. Is this valid?
\begin{enumerate}
\item Yes, this is standard and correct.
\item It can be valid but has a subtle issue: feature selection using the full data introduces selection bias. Ideally, LASSO selection should be done inside each cross-validation fold to avoid leakage.
\item No, LASSO should not be used for feature selection.
\item No, logistic regression should be used directly.
\end{enumerate} \end{block}
\begin{exampleblock}{Answer: B}\end{exampleblock}
\begin{block}{Explanation}
Selection on the full data leaks information. Nested CV avoids this by doing selection within each fold. Options A, C, and D misstate the issue or the role of LASSO.
\end{block}
\end{frame}

% ---------------- MODULE 2 — TOOLS AND TECHNIQUES ----------------

\begin{frame}{Q271 --- K-Means Feature Scaling}
\begin{block}{Question}
A K-means clustering is run on customer data with features: annual spend (USD 100--1M), tenure (months, 1--240), and number of products (1--10). The result is dominated by annual spend. What should the analyst do?
\end{block}
\begin{enumerate}
\item Remove annual spend.
\item Standardize or normalize the features before running K-means.
\item Use fewer observations.
\item Increase K.
\end{enumerate}
\begin{exampleblock}{Answer: B}\end{exampleblock}
\begin{block}{Explanation}
K-means uses distances; features on different scales dominate. Scaling fixes this. Options A, C, and D do not address the scale issue.
\end{block}
\end{frame}

\begin{frame}{Q272 --- Silhouette Score Interpretation}
\begin{block}{Question}
A clustering with K=4 has average silhouette score 0.32. With K=5, the score is 0.15. What do these scores indicate?
\begin{enumerate}
\item K=5 is better because the score is lower.
\item K=4 is preferred because the average silhouette score is higher, indicating better-defined clusters.
\item Neither is good.
\item Silhouette scores only apply to K=2.
\end{enumerate}  \end{block}
\begin{exampleblock}{Answer: B}\end{exampleblock}
\begin{block}{Explanation}
Higher silhouette = better-separated clusters. K=4 is preferred. Options A, C, and D misstate the interpretation.
\end{block}
\end{frame}

\begin{frame}{Q273 --- Dendrogram Cut}
\begin{block}{Question}
A dendrogram on 500 customers shows two very long vertical branches near the top, then progressively shorter merges. How many clusters are suggested?
\end{block}
\begin{enumerate}
\item 500 (each customer is its own cluster).
\item 2 (the two branches separated by the long vertical lines).
\item 10 (based on the number of intermediate branches).
\item 1 (all merged).
\end{enumerate}
\begin{exampleblock}{Answer: B}\end{exampleblock}
\begin{block}{Explanation}
Long vertical lines indicate significant separation. Cutting above the long vertical lines yields 2 clusters. Options A, C, and D misread the dendrogram.
\end{block}
\end{frame}

\begin{frame}{Q274 --- Random Forest vs Gradient Boosting}
\begin{block}{Question}
A quant team is deciding between random forest and gradient-boosted trees for a medium-sized tabular classification problem (10,000 rows, 40 features). Which is generally true?
\end{block}
\begin{enumerate}
\item Random forest is always better.
\item Both are competitive; gradient boosting often achieves slightly higher accuracy with careful tuning, while random forest is generally more robust to hyperparameter choices and slightly faster to train.
\item Gradient boosting is always better.
\item Neither is suitable.
\end{enumerate}
\begin{exampleblock}{Answer: B}\end{exampleblock}
\begin{block}{Explanation}
Both are strong; boosting often wins on accuracy with tuning, while RF is easier to use. Options A, C, and D are overstated.
\end{block}
\end{frame}

\begin{frame}{Q275 --- SVM and Kernel Choice}
\begin{block}{Question}
An SVM with a linear kernel underperforms; with an RBF kernel it performs better. What does this suggest about the data?
\end{block}
\begin{enumerate}
\item The classes are linearly separable.
\item The classes are not linearly separable, and a non-linear decision boundary is required.
\item The data is noisy.
\item The SVM is broken.
\end{enumerate}
\begin{exampleblock}{Answer: B}\end{exampleblock}
\begin{block}{Explanation}
The improved performance with a non-linear kernel indicates the boundary is non-linear. Options A, C, and D are not supported.
\end{block}
\end{frame}

\begin{frame}{Q276 --- KNN in High Dimensions}
\begin{block}{Question}
KNN performs poorly on a dataset with 200 features and 1,000 observations. What is a likely cause and remedy?
\begin{enumerate} 
\item Too many observations; reduce.
\item Curse of dimensionality; use dimensionality reduction (PCA, feature selection) or a different algorithm.
\item Wrong distance metric; use Manhattan.
\item Increase k.
\end{enumerate}\end{block}
\begin{exampleblock}{Answer: B}\end{exampleblock}
\begin{block}{Explanation}
High-dimensional spaces suffer from distance concentration, making KNN unreliable. Dimensionality reduction or alternative algorithms are recommended. Options A, C, and D are secondary.
\end{block}
\end{frame}

\begin{frame}{Q277 --- Decision Trees and Pruning}
\begin{block}{Question}
A fully grown decision tree achieves 100\% accuracy on training data and poor accuracy on test data. What should be done?
\begin{enumerate}
\item Grow the tree deeper.
\item Prune the tree (or apply pre-pruning constraints on depth, minimum samples per leaf, etc.) to reduce overfitting.
\item Use a different dataset.
\item Remove all features.
\end{enumerate} \end{block}
\begin{exampleblock}{Answer: B}\end{exampleblock}
\begin{block}{Explanation}
Unpruned trees overfit; pruning reduces variance. Options A, C, and D do not address overfitting.
\end{block}
\end{frame}

\begin{frame}{Q278 --- Ensemble Diversity}
\begin{block}{Question}
An ensemble of 50 identical decision trees performs no better than a single tree. What is the explanation?
\end{block}
\begin{enumerate}
\item Ensembles always improve performance.
\item Diversity is essential: identical trees give identical predictions, so the ensemble has no variance reduction benefit. Bagging and feature randomness introduce the needed diversity.
\item The trees are too shallow.
\item The ensemble is too small.
\end{enumerate}
\begin{exampleblock}{Answer: B}\end{exampleblock}
\begin{block}{Explanation}
Ensembles improve performance when base learners are diverse and individually reasonable. Identical trees defeat the purpose. Options A, C, and D miss the point.
\end{block}
\end{frame}

\begin{frame}{Q279 --- Naive Bayes Assumptions}
\begin{block}{Question}
A Naive Bayes classifier performs surprisingly well on text despite the conditional independence assumption being violated. Why?
\begin{enumerate}
\item The assumption is satisfied in text data.
\item Classification depends only on ranking class posteriors, not on exact probability calibration; Naive Bayes often ranks correctly even when probabilities are miscalibrated.
\item Text data has no dependencies.
\item Naive Bayes ignores the assumption in practice.
\end{enumerate}\end{block}
\begin{exampleblock}{Answer: B}\end{exampleblock}
\begin{block}{Explanation}
Naive Bayes often ranks classes correctly even with miscalibrated probabilities because the decision boundary depends on the maximum, not on calibration. Options A, C, and D are wrong.
\end{block}
\end{frame}

\begin{frame}{Q280 --- Logistic Regression vs LDA}
\begin{block}{Question}
A binary classification problem has features that are approximately normally distributed within each class with equal covariance matrices. Which is theoretically preferable: logistic regression or linear discriminant analysis?
\end{block}
\begin{enumerate}
\item Logistic regression, always.
\item LDA, because when its assumptions (normal features, equal covariance) hold, it is the optimal Bayes classifier.
\item Both are equally good regardless of assumptions.
\item SVM.
\end{enumerate}
\begin{exampleblock}{Answer: B}\end{exampleblock}
\begin{block}{Explanation}
Under normality and equal covariance, LDA is the Bayes-optimal classifier. Logistic regression is more robust to assumption violations and often used when these don't hold. Options A, C, and D are not correct.
\end{block}
\end{frame}

\begin{frame}{Q281 --- Regression vs Classification}
\begin{block}{Question}
An analyst wants to predict whether a loan will default or not. Which type of model is appropriate?
\begin{enumerate}
\item Linear regression.
\item Logistic regression (or another classification model).
\item Polynomial regression.
\item Ridge regression.
\end{enumerate}\end{block}
\begin{exampleblock}{Answer: B}\end{exampleblock}
\begin{block}{Explanation}
Default is a binary outcome; classification models are appropriate. Linear and polynomial regressions predict continuous outputs; Ridge is regularized linear regression. Option B is correct.
\end{block}
\end{frame}

\begin{frame}{Q282 --- LDA Assumptions}
\begin{block}{Question}
Which of the following is NOT an assumption of standard Linear Discriminant Analysis?
\begin{enumerate}
\item Features are normally distributed within each class.
\item Classes share a common covariance matrix.
\item Features are continuous.
\item Classes have equal prior probabilities.
\end{enumerate}\end{block}
\begin{exampleblock}{Answer: D}\end{exampleblock}
\begin{block}{Explanation}
LDA allows for unequal priors (estimated from data or specified). The other three assumptions are standard. Option D is not required.
\end{block}
\end{frame}

\begin{frame}{Q283 --- ROC AUC in Imbalanced Data}
\begin{block}{Question}
Two models are compared on a fraud dataset with 1\% positives. Model A has AUC 0.90; Model B has AUC 0.75. But Model B has a higher precision-recall AUC (PR-AUC). Which model is better?
\end{block}
\begin{enumerate}
\item Model A, because ROC AUC is always the criterion.
\item Depends on the use case; PR-AUC is often more informative for highly imbalanced data, so Model B might be preferable if precision at high recall matters.
\item Model B, always.
\item Neither; discard both.
\end{enumerate}
\begin{exampleblock}{Answer: B}\end{exampleblock}
\begin{block}{Explanation}
For imbalanced data, PR-AUC often better reflects performance on the minority class. Choice depends on business context. Options A, C, and D are overstated.
\end{block}
\end{frame}

\begin{frame}{Q284 --- Classification Threshold and Business Cost}
\begin{block}{Question}
A bank has a default prediction model. Cost of approving a defaulter (FP) is \$50k; cost of rejecting a good borrower (FN) is \$5k. How should the threshold be chosen?
\begin{enumerate}
\item Threshold = 0.5 always.
\item Threshold should be higher than 0.5 to reduce false positives, since FP cost >> FN cost.
\item Threshold should be lower than 0.5 to increase recall.
\item Threshold choice is irrelevant.
\end{enumerate}\end{block}
\begin{exampleblock}{Answer: B}\end{exampleblock}
\begin{block}{Explanation}
When FP is much costlier than FN, raising the threshold reduces FPs. Options A, C, and D misapply the cost matrix.
\end{block}
\end{frame}

\begin{frame}{Q285 --- Precision, Recall, and F1}
\begin{block}{Question}
Model A has precision 0.9, recall 0.2. Model B has precision 0.4, recall 0.9. Which model is likely appropriate for a fraud-detection system where missing fraud is very costly?
\end{block}
\begin{enumerate}
\item Model A, because precision is high.
\item Model B, because recall is high, catching more actual frauds.
\item Both are equivalent.
\item Neither.
\end{enumerate}
\begin{exampleblock}{Answer: B}\end{exampleblock}
\begin{block}{Explanation}
When missing positives is very costly, recall is prioritised. Model B's high recall is preferable. Options A, C, and D are less appropriate.
\end{block}
\end{frame}

\begin{frame}{Q286 --- Confusion Matrix Arithmetic}
\begin{block}{Question}
A confusion matrix shows TP=80, FP=20, FN=40, TN=860. Compute the F1 score.
\end{block}
\begin{enumerate}
\item 0.615.
\item 0.727.
\item 0.842.
\item 0.500.
\end{enumerate} 
\begin{exampleblock}{Answer: B}\end{exampleblock}
\begin{block}{Explanation}
Precision  0.80; Options A, C, and D miscompute.
\end{block}
\end{frame}

\begin{frame}{Q287 --- Regularization Choice}
\begin{block}{Question}
A team has 50 features and 200 observations. They expect some features to be irrelevant. Which regularization is most appropriate?
\begin{enumerate}
\item L2 (Ridge).
\item L1 (LASSO).
\item Elastic Net.
\item None.
\end{enumerate} \end{block}
\begin{exampleblock}{Answer: B}\end{exampleblock}
\begin{block}{Explanation}
When many features may be irrelevant, LASSO's feature selection is valuable. Elastic Net is also reasonable if features are correlated. Options A and D are less appropriate.
\end{block}
\end{frame}

\begin{frame}{Q288 --- Cross-Validation in Small Samples}
\begin{block}{Question}
With only 100 observations, why might leave-one-out cross-validation (LOOCV) be preferred over k-fold?
\begin{enumerate}
\item Because it is faster.
\item Because it uses almost all data for training in each fold, reducing bias in the performance estimate when the sample is very small.
\item Because it is more random.
\item Because it avoids overfitting.
\end{enumerate} \end{block}
\begin{exampleblock}{Answer: B}\end{exampleblock}
\begin{block}{Explanation}
LOOCV maximises training data per fold; useful for very small samples. It is computationally expensive and higher variance, but lower bias. Options A, C, and D misstate the trade-off.
\end{block}
\end{frame}

\begin{frame}{Q289 --- Bootstrap vs Cross-Validation}
\begin{block}{Question}
Which is a key difference between bootstrapping and k-fold cross-validation for model evaluation?
\begin{enumerate}
\item Bootstrapping samples with replacement; k-fold partitions the data without replacement.
\item Bootstrapping is always better.
\item k-fold can only be used for classification.
\item They are identical.
\end{enumerate} \end{block}
\begin{exampleblock}{Answer: A}\end{exampleblock}
\begin{block}{Explanation}
Bootstrapping uses sampling with replacement; k-fold partitions the data. Both estimate generalisation. Options B, C, and D misstate the distinction.
\end{block}
\end{frame}

\begin{frame}{Q290 --- Regularization and Feature Scaling}
\begin{block}{Question}
Why is standardization recommended before applying L1 or L2 regularization?
\begin{enumerate}
\item It is not recommended.
\item Because regularization penalises coefficients, and coefficients on differently scaled features are not directly comparable; standardizing puts them on equal footing.
\item It removes outliers.
\item It speeds up training.
\end{enumerate} \end{block}
\begin{exampleblock}{Answer: B}\end{exampleblock}
\begin{block}{Explanation}
Regularization penalises coefficient magnitude. If features are not on comparable scales, the penalty is unfairly applied. Options A, C, and D are not the reason.
\end{block}
\end{frame}

\begin{frame}{Q291 --- Neural Network Activation Functions}
\begin{block}{Question}
A hidden layer uses sigmoid activation. After many layers, gradients vanish and training stalls. Which alternative activation is most likely to help?
\begin{enumerate}
\item Another sigmoid layer.
\item ReLU (or its variants like Leaky ReLU, ELU).
\item Softmax.
\item Linear.
\end{enumerate} \end{block}
\begin{exampleblock}{Answer: B}\end{exampleblock}
\begin{block}{Explanation}
ReLU avoids saturation for positive inputs and mitigates vanishing gradients. Softmax is for output; linear is not typically used in hidden layers. Options A, C, and D are less appropriate.
\end{block}
\end{frame}

\begin{frame}{Q292 --- Weight Initialization}
\begin{block}{Question}
Why is Xavier or He initialization used for neural network weights?
\begin{enumerate}
\item To initialize weights to zero.
\item To maintain variance of activations and gradients across layers, avoiding vanishing/exploding gradients.
\item To speed up training at the cost of accuracy.
\item To reduce model size.
\end{enumerate} \end{block}
\begin{exampleblock}{Answer: B}\end{exampleblock}
\begin{block}{Explanation}
Xavier/He initialization scales weights based on layer sizes, preserving signal magnitude. Options A, C, and D are incorrect.
\end{block}
\end{frame}

\begin{frame}{Q293 --- Batch Normalization}
\begin{block}{Question}
Batch normalization improves deep network training primarily by:
\begin{enumerate}
\item Reducing model complexity.
\item Normalizing layer inputs, which stabilises training, mitigates internal covariate shift, and often allows higher learning rates.
\item Eliminating the need for activation functions.
\item Reducing the number of parameters.
\end{enumerate} \end{block}
\begin{exampleblock}{Answer: B}\end{exampleblock}
\begin{block}{Explanation}
Batch norm normalizes intermediate activations, improving stability and convergence. Options A, C, and D are incorrect.
\end{block}
\end{frame}

\begin{frame}{Q294 --- Dropout at Inference}
\begin{block}{Question}
Dropout is used during training. What happens at inference time?
\begin{enumerate}
\item Dropout is applied again.
\item Dropout is turned off; weights are scaled to account for the expected effect of dropout during training.
\item Model is retrained without dropout.
\item Model is discarded.
\end{enumerate}\end{block}
\begin{exampleblock}{Answer: B}\end{exampleblock}
\begin{block}{Explanation}
Dropout is a training-time regularizer. At inference, dropout is disabled and weights are scaled (or the framework handles it). Options A, C, and D are incorrect.
\end{block}
\end{frame}

\begin{frame}{Q295 --- Early Stopping}
\begin{block}{Question}
Early stopping in neural network training means:
\begin{enumerate}
\item Stopping before training starts.
\item Monitoring validation error during training and stopping when it stops improving, preventing overfitting.
\item Stopping after a fixed number of epochs regardless of performance.
\item Stopping when training error is zero.
\end{enumerate} \end{block}
\begin{exampleblock}{Answer: B}\end{exampleblock}
\begin{block}{Explanation}
Early stopping is a form of regularization: training halts when validation performance degrades. Options A, C, and D are different or incorrect.
\end{block}
\end{frame}

\begin{frame}{Q296 --- Convolutional Layers}
\begin{block}{Question}
Why do CNNs use shared weights (convolutional filters) across spatial positions?
\begin{enumerate}
\item To reduce training time only.
\item To exploit spatial invariance and dramatically reduce the number of parameters compared to fully connected layers.
\item Because convolution is faster than matrix multiplication.
\item To avoid using activation functions.
\end{enumerate}\end{block}
\begin{exampleblock}{Answer: B}\end{exampleblock}
\begin{block}{Explanation}
Weight sharing reduces parameters and provides translational invariance. Options A, C, and D miss the architectural rationale.
\end{block}
\end{frame}

\begin{frame}{Q297 --- Pooling Layers}
\begin{block}{Question}
What is the primary purpose of pooling layers in CNNs?
\begin{enumerate}
\item To increase spatial resolution.
\item To reduce spatial dimensions (downsampling), control overfitting, and provide translation invariance.
\item To add parameters.
\item To replace convolution.
\end{enumerate} \end{block}
\begin{exampleblock}{Answer: B}\end{exampleblock}
\begin{block}{Explanation}
Pooling reduces spatial size and provides robustness to small translations. Options A, C, and D are incorrect.
\end{block}
\end{frame}

\begin{frame}{Q298 --- Attention Mechanism}
\begin{block}{Question}
In a transformer, self-attention computes pairwise interactions between tokens. What is the primary benefit?
\begin{enumerate}
\item Faster training on small data.
\item Modelling long-range dependencies and contextual relationships regardless of token distance.
\item Reducing parameter count.
\item Handling images only.
\end{enumerate}\end{block}
\begin{exampleblock}{Answer: B}\end{exampleblock}
\begin{block}{Explanation}
Self-attention captures relationships between tokens regardless of distance, unlike RNNs. Options A, C, and D are incorrect.
\end{block}
\end{frame}

\begin{frame}{Q299 --- Positional Encoding}
\begin{block}{Question}
Why do transformers need positional encodings?
\begin{enumerate} 
\item Because transformers lack recurrence and therefore have no inherent notion of order.
\item Because transformers are too slow without them.
\item Because they help with image processing.
\item Because they increase model size.
\end{enumerate} \end{block}
\begin{exampleblock}{Answer: A}\end{exampleblock}
\begin{block}{Explanation}
Self-attention is permutation-invariant; positional encodings inject order information. Options B, C, and D are incorrect.
\end{block}
\end{frame}

\begin{frame}{Q300 --- Transformers vs RNNs}
\begin{block}{Question}
A key advantage of transformers over RNNs is:
\begin{enumerate}
\item Less computational cost.
\item Parallelism over sequence positions and better handling of long-range dependencies.
\item Better handling of small datasets.
\item Simpler architecture.
\end{enumerate}\end{block}
\begin{exampleblock}{Answer: B}\end{exampleblock}
\begin{block}{Explanation}
Transformers parallelise across tokens and capture long-range dependencies via attention. RNNs are sequential. Options A, C, and D are less accurate.
\end{block}
\end{frame}

% ---------------- MODULE 3 — RISK AND RISK FACTORS ----------------

\begin{frame}{Q301 --- Fairness Metrics}
\begin{block}{Question}
Two groups have equal true positive rates but unequal false positive rates under a model. Which fairness metric does the model satisfy?
\begin{enumerate}
\item Demographic parity.
\item Equal opportunity.
\item Predictive rate parity.
\item None.
\end{enumerate} \end{block}
\begin{exampleblock}{Answer: B}\end{exampleblock}
\begin{block}{Explanation}
Equal opportunity requires equal TPR across groups. The model satisfies that but violates equalized odds (which requires equal FPR too) and demographic parity. Options A, C, and D are not satisfied.
\end{block}
\end{frame}

\begin{frame}{Q302 --- Bias Types}
\begin{block}{Question}
An algorithm trained to predict job performance is disproportionately rejecting female candidates. Removing gender from features does not fix the issue. What is the most likely cause?
\begin{enumerate}
\item Random noise.
\item Proxy discrimination via correlated features (e.g., university, career gaps).
\item The training set is too small.
\item The target is wrong.
\end{enumerate}\end{block}
\begin{exampleblock}{Answer: B}\end{exampleblock}
\begin{block}{Explanation}
Protected attributes are often captured by correlated proxies. Options A, C, and D are secondary or incorrect.
\end{block}
\end{frame}

\begin{frame}{Q303 --- Fairness Impossibility}
\begin{block}{Question}
Which of the following cannot be simultaneously satisfied when base rates differ between groups?
\begin{enumerate}
\item Accuracy and precision.
\item Demographic parity, predictive rate parity, and equal opportunity.
\item Recall and F1.
\item ROC and AUC.
\end{enumerate}\end{block}
\begin{exampleblock}{Answer: B}\end{exampleblock}
\begin{block}{Explanation}
The Kleinberg-Chouldechova impossibility result applies to these three fairness metrics. Options A, C, and D are not the impossibility case.
\end{block}
\end{frame}

\begin{frame}{Q304 --- Historical Bias}
\begin{block}{Question}
A model trained on 20 years of hiring data favours male candidates. Which type of bias?
\begin{enumerate}
\item Measurement bias.
\item Historical bias.
\item Sampling bias.
\item Deployment bias.
\end{enumerate}\end{block}
\begin{exampleblock}{Answer: B}\end{exampleblock}
\begin{block}{Explanation}
Historical bias reflects past discrimination encoded in data. Options A, C, and D are other sources.
\end{block}
\end{frame}

\begin{frame}{Q305 --- Feedback Loops}
\begin{block}{Question}
A model trained on data generated by its own past predictions can become increasingly extreme. What is this phenomenon?
\begin{enumerate}
\item Random noise.
\item Feedback loop or self-reinforcing bias.
\item Overfitting.
\item Underfitting.
\end{enumerate}\end{block}
\begin{exampleblock}{Answer: B}\end{exampleblock}
\begin{block}{Explanation}
Feedback loops occur when model outputs influence future training data, amplifying biases. Options A, C, and D are different phenomena.
\end{block}
\end{frame}

\begin{frame}{Q306 --- Interpretability vs Explainability}
\begin{block}{Question}
A neural network's internal weights are not human-interpretable, but SHAP values explain individual predictions. What is this model?
\begin{enumerate}
\item Interpretable and explainable.
\item Explainable but not interpretable.
\item Interpretable but not explainable.
\item Neither.
\end{enumerate}\end{block}
\begin{exampleblock}{Answer: B}\end{exampleblock}
\begin{block}{Explanation}
SHAP provides post-hoc explainability; the network's internals remain opaque. Options A, C, and D misstate.
\end{block}
\end{frame}

\begin{frame}{Q307 --- Transparency}
	\begin{block}{Question}
		Which is an example of opacity by ``specialised knowledge''?
		\begin{enumerate}
			\item Being unaware a decision was made by an algorithm.
			\item Even with access to the code, only experts can interpret it.
			\item Hiding the algorithm.
			\item Using encrypted data.
	\end{enumerate}  \end{block}
	\begin{exampleblock}{Answer: B}\end{exampleblock}
	\begin{block}{Explanation}
		Specialised-knowledge opacity arises when understanding requires expertise. Option A is secrecy; C is confidentiality; D is unrelated.
	\end{block}
\end{frame}

\begin{frame}{Q308 --- Automation Bias}
	\begin{block}{Question}
		An operator accepts AI recommendations without independent verification. What is this?
		\begin{enumerate}
			\item Confirmation bias.
			\item Automation bias.
			\item Selection bias.
			\item Survivorship bias.
	\end{enumerate}  \end{block}
	\begin{exampleblock}{Answer: B}\end{exampleblock}
	\begin{block}{Explanation}
		Automation bias is over-reliance on automated systems. Options A, C, and D are different biases.
	\end{block}
\end{frame}

\begin{frame}{Q309 --- Reputational Risk Sources}
	\begin{block}{Question}
		Which of the following are recognised sources of AI-related reputational risk?
		\begin{enumerate}
			\item Privacy breaches, algorithmic bias, and lack of explainability.
			\item Only privacy.
			\item Only financial losses.
			\item Only regulatory fines.
	\end{enumerate}  \end{block}
	\begin{exampleblock}{Answer: A}\end{exampleblock}
	\begin{block}{Explanation}
		Privacy, bias, and explainability are the three primary sources of AI-related reputational risk. Options B, C, and D are too narrow.
	\end{block}
\end{frame}

\begin{frame}{Q310 --- Societal Risks of GenAI}
	\begin{block}{Question}
		Which is a societal risk of GenAI?
		\begin{enumerate}
			\item High compute cost.
			\item Misinformation at scale and deepfake-enabled fraud.
			\item Slow inference.
			\item Privacy only.
	\end{enumerate}  \end{block}
	\begin{exampleblock}{Answer: B}\end{exampleblock}
	\begin{block}{Explanation}
		GenAI's societal risks include misinformation and deepfakes. Options A and C are technical; D is too narrow.
	\end{block}
\end{frame}

\begin{frame}{Q311 --- Autonomy Violation}
	\begin{block}{Question}
		An AI system overrides a patient's stated preference for a treatment. Which principle is violated?
		\begin{enumerate}
			\item Nonmaleficence.
			\item Beneficence.
			\item Autonomy.
			\item Explainability.
	\end{enumerate}  \end{block}
	\begin{exampleblock}{Answer: C}\end{exampleblock}
	\begin{block}{Explanation}
		Autonomy is the patient's right to decide. Options A, B, and D are distinct principles.
	\end{block}
\end{frame}

\begin{frame}{Q312 --- Nonmaleficence vs Beneficence}
	\begin{block}{Question}
		Which is the difference between nonmaleficence and beneficence?
		\begin{enumerate}
			\item Nonmaleficence: avoid harm; beneficence: actively do good.
			\item Both mean the same.
			\item Beneficence: avoid harm; nonmaleficence: actively do good.
			\item Neither applies to AI.
	\end{enumerate}  \end{block}
	\begin{exampleblock}{Answer: A}\end{exampleblock}
	\begin{block}{Explanation}
		Nonmaleficence = do no harm; beneficence = do good. Options B, C, and D misstate.
	\end{block}
\end{frame}

\begin{frame}{Q313 --- Justice Conflicts}
	\begin{block}{Question}
		Egalitarianism and meritocracy can conflict. What does this illustrate about justice in AI ethics?
		\begin{enumerate}
			\item Justice is uncontroversial.
			\item Different distributive justice concepts can conflict, requiring trade-offs.
			\item Justice is a mathematical rule.
			\item Justice is irrelevant to AI.
	\end{enumerate}  \end{block}
	\begin{exampleblock}{Answer: B}\end{exampleblock}
	\begin{block}{Explanation}
		Justice concepts can conflict; choosing between them is a normative judgment. Options A, C, and D are incorrect.
	\end{block}
\end{frame}

\begin{frame}{Q314 --- Human Autonomy in Practice}
	\begin{block}{Question}
		A robo-advisor nudges users toward certain products through manipulative design. Which principle does this threaten?
		\begin{enumerate}
			\item Autonomy.
			\item Justice.
			\item Nonmaleficence.
			\item Explainability.
	\end{enumerate}  \end{block}
	\begin{exampleblock}{Answer: A}\end{exampleblock}
	\begin{block}{Explanation}
		Manipulative design undermines autonomy. Options B, C, and D are different principles.
	\end{block}
\end{frame}

\begin{frame}{Q315 --- Manipulation and AI}
	\begin{block}{Question}
		Which of the following is an example of AI-driven manipulation?
		\begin{enumerate}
			\item A calculator computes sums.
			\item A recommendation engine leverages detailed user profiles to shape opinions or behaviour without user awareness.
			\item A chatbot answers FAQs.
			\item An OCR system reads text.
	\end{enumerate}  \end{block}
	\begin{exampleblock}{Answer: B}\end{exampleblock}
	\begin{block}{Explanation}
		Manipulation involves leveraging profiling to influence users. Options A, C, and D are not manipulative.
	\end{block}
\end{frame}

\begin{frame}{Q316 --- Existential Risk}
	\begin{block}{Question}
		Which of the following is a key concern in AI existential risk discussions?
		\begin{enumerate}
			\item Model accuracy.
			\item Superintelligence and the alignment problem.
			\item Data storage costs.
			\item IT infrastructure.
	\end{enumerate}  \end{block}
	\begin{exampleblock}{Answer: B}\end{exampleblock}
	\begin{block}{Explanation}
		Existential risk focuses on superintelligence and alignment with human values. Options A, C, and D are operational.
	\end{block}
\end{frame}

\begin{frame}{Q317 --- Global Inequality}
	\begin{block}{Question}
		What is a concern regarding AI and global inequality?
		\begin{enumerate}
			\item AI is equally accessible everywhere.
			\item Countries with advanced AI capabilities may pull ahead economically, widening global inequality.
			\item AI reduces inequality.
			\item AI is irrelevant to inequality.
	\end{enumerate}  \end{block}
	\begin{exampleblock}{Answer: B}\end{exampleblock}
	\begin{block}{Explanation}
		Concentration of AI capability can widen global economic and political gaps. Options A, C, and D are incorrect or unsupported.
	\end{block}
\end{frame}

\begin{frame}{Q318 --- Misinformation Campaigns}
	\begin{block}{Question}
		Why is AI particularly relevant to misinformation campaigns?
		\begin{enumerate}
			\item AI cannot generate text.
			\item AI can generate realistic content at scale cheaply, amplifying misinformation and disinformation.
			\item AI only produces accurate information.
			\item AI is too slow to matter.
	\end{enumerate}  \end{block}
	\begin{exampleblock}{Answer: B}\end{exampleblock}
	\begin{block}{Explanation}
		Generative AI makes content creation cheap and scalable, increasing misinformation risk. Options A, C, and D are incorrect.
	\end{block}
\end{frame}

\begin{frame}{Q319 --- Privacy and Surveillance}
	\begin{block}{Question}
		Which is a privacy concern associated with AI?
		\begin{enumerate}
			\item AI improves accuracy only.
			\item AI systems can infer sensitive attributes or habits from seemingly innocuous data, potentially enabling surveillance.
			\item AI always protects privacy.
			\item AI has no privacy implications.
	\end{enumerate}  \end{block}
	\begin{exampleblock}{Answer: B}\end{exampleblock}
	\begin{block}{Explanation}
		AI can infer sensitive information from data, raising privacy and surveillance risks. Options A, C, and D are incorrect.
	\end{block}
\end{frame}

\begin{frame}{Q320 --- GenAI Functional Risks}
	\begin{block}{Question}
		Which of the following is a functional risk of GenAI?
		\begin{enumerate}
			\item Hallucination of plausible but incorrect outputs.
			\item Reduced model size.
			\item Faster inference.
			\item Lower training cost.
	\end{enumerate}  \end{block}
	\begin{exampleblock}{Answer: A}\end{exampleblock}
	\begin{block}{Explanation}
		Hallucination is a functional risk of GenAI. Options B, C, and D are benefits, not risks.
	\end{block}
\end{frame}

% ---------------- MODULE 4 — RESPONSIBLE AND ETHICAL AI ----------------

\begin{frame}{Q321 --- Utilitarian Trade-offs}
	\begin{block}{Question}
		A self-driving car algorithm is tuned to minimise total harm, potentially at the cost of the passenger. Which ethical framework best supports this?
		\begin{enumerate}
			\item Deontology.
			\item Virtue ethics.
			\item Utilitarianism.
			\item Divine command theory.
	\end{enumerate}  \end{block}
	\begin{exampleblock}{Answer: C}\end{exampleblock}
	\begin{block}{Explanation}
		Utilitarianism maximises aggregate welfare, potentially overriding individual rights. Options A, B, and D are different frameworks.
	\end{block}
\end{frame}

\begin{frame}{Q322 --- Deontological Constraints}
	\begin{block}{Question}
		A company refuses to use AI for certain tasks even if it would improve aggregate outcomes, because the tasks violate fundamental rights. Which ethical framework guides this?
		\begin{enumerate}
			\item Utilitarianism.
			\item Deontology.
			\item Virtue ethics.
			\item Pragmatism.
	\end{enumerate}  \end{block}
	\begin{exampleblock}{Answer: B}\end{exampleblock}
	\begin{block}{Explanation}
		Deontology prioritises duties and rights over consequences. Options A, C, and D are different.
	\end{block}
\end{frame}

\begin{frame}{Q323 --- Virtue Ethics in AI}
	\begin{block}{Question}
		What question would a virtue ethicist ask when designing an AI system?
		\begin{enumerate}
			\item What maximises utility?
			\item What would a virtuous person/organisation do?
			\item What does the law say?
			\item What is cheapest?
	\end{enumerate}  \end{block}
	\begin{exampleblock}{Answer: B}\end{exampleblock}
	\begin{block}{Explanation}
		Virtue ethics focuses on character and virtuous behaviour. Options A, C, and D are other frameworks or criteria.
	\end{block}
\end{frame}

\begin{frame}{Q324 --- Consent and Privacy}
	\begin{block}{Question}
		Which of the following is essential for ethical data collection under GDPR?
		\begin{enumerate}
			\item No notice.
			\item Explicit consent (in many cases) or another lawful basis, data minimisation, and transparent purpose specification.
			\item Unlimited retention.
			\item Sharing with third parties.
	\end{enumerate}  \end{block}
	\begin{exampleblock}{Answer: B}\end{exampleblock}
	\begin{block}{Explanation}
		GDPR requires lawful basis, minimisation, and transparency. Options A, C, and D violate GDPR.
	\end{block}
\end{frame}

\begin{frame}{Q325 --- Right to be Forgotten}
	\begin{block}{Question}
		Which of the following best captures the GDPR right to be forgotten?
		\begin{enumerate}
			\item Right to access data.
			\item Right to request erasure of personal data under certain conditions.
			\item Right to sell data.
			\item Right to receive compensation for processing.
	\end{enumerate}  \end{block}
	\begin{exampleblock}{Answer: B}\end{exampleblock}
	\begin{block}{Explanation}
		The right to erasure (Art. 17) allows deletion under specified conditions. Options A, C, and D are other rights.
	\end{block}
\end{frame}

\begin{frame}{Q326 --- Contextual Integrity}
	\begin{block}{Question}
		What does contextual integrity require?
		\begin{enumerate}
			\item Data can be used freely once collected.
			\item Data use must respect the norms of the context in which it was provided.
			\item Data should be encrypted.
			\item Data should be sold.
	\end{enumerate}  \end{block}
	\begin{exampleblock}{Answer: B}\end{exampleblock}
	\begin{block}{Explanation}
		Contextual integrity (Nissenbaum) requires data use to align with context-specific norms. Options A, C, and D misstate or miss the concept.
	\end{block}
\end{frame}

\begin{frame}{Q327 --- EU AI Act}
	\begin{block}{Question}
		Which of the following is explicitly prohibited by the EU AI Act?
		\begin{enumerate}
			\item Fraud detection.
			\item Social scoring and manipulative AI exploiting vulnerabilities.
			\item Credit scoring.
			\item Customer service chatbots.
	\end{enumerate}  \end{block}
	\begin{exampleblock}{Answer: B}\end{exampleblock}
	\begin{block}{Explanation}
		Social scoring, manipulative AI, and exploitation of vulnerabilities are prohibited. Other uses are regulated but permitted.
	\end{block}
\end{frame}

\begin{frame}{Q328 --- GDPR Extraterritoriality}
	\begin{block}{Question}
		A US firm serves EU customers. Does GDPR apply?
		\begin{enumerate}
			\item No.
			\item Yes, GDPR applies extraterritorially when serving EU data subjects.
			\item Only if the firm has an EU office.
			\item Only if the customer complains.
	\end{enumerate}  \end{block}
	\begin{exampleblock}{Answer: B}\end{exampleblock}
	\begin{block}{Explanation}
		GDPR Art. 3 provides extraterritorial scope. Options A, C, and D are incorrect.
	\end{block}
\end{frame}

\begin{frame}{Q329 --- NIST AI RMF}
	\begin{block}{Question}
		The NIST AI RMF functions are:
		\begin{enumerate}
			\item Train, Test, Deploy, Monitor.
			\item Govern, Map, Measure, Manage.
			\item Design, Build, Audit, Retire.
			\item Plan, Do, Check, Act.
	\end{enumerate}  \end{block}
	\begin{exampleblock}{Answer: B}\end{exampleblock}
	\begin{block}{Explanation}
		NIST AI RMF functions: Govern, Map, Measure, Manage. Options A, C, and D are other frameworks.
	\end{block}
\end{frame}

\begin{frame}{Q330 --- Kantian Ethics}
	\begin{block}{Question}
		Kant's categorical imperative asks us to:
		\begin{enumerate}
			\item Maximise utility.
			\item Act only on maxims that could be universal laws.
			\item Cultivate virtue.
			\item Follow the sovereign.
	\end{enumerate}  \end{block}
	\begin{exampleblock}{Answer: B}\end{exampleblock}
	\begin{block}{Explanation}
		Kant's imperative demands universalizability. Options A, C, and D are other traditions.
	\end{block}
\end{frame}

\begin{frame}{Q331 --- Ethical Frameworks Comparison}
	\begin{block}{Question}
		Which statement correctly distinguishes the three main ethical frameworks?
		\begin{enumerate}
			\item All three are identical.
			\item Consequentialism focuses on outcomes; deontology on duties; virtue ethics on character.
			\item All three focus on outcomes.
			\item All three ignore outcomes.
	\end{enumerate}  \end{block}
	\begin{exampleblock}{Answer: B}\end{exampleblock}
	\begin{block}{Explanation}
		The three frameworks differ in focus: consequences, duties, character. Options A, C, and D are incorrect.
	\end{block}
\end{frame}

\begin{frame}{Q332 --- Privacy Principles}
	\begin{block}{Question}
		Which of the following is NOT typically considered a privacy best practice?
		\begin{enumerate}
			\item Data minimisation.
			\item Right to be forgotten.
			\item Contextual integrity.
			\item Unlimited data retention.
	\end{enumerate}  \end{block}
	\begin{exampleblock}{Answer: D}\end{exampleblock}
	\begin{block}{Explanation}
		Unlimited retention is contrary to privacy best practices. Options A, B, and C are established practices.
	\end{block}
\end{frame}

\begin{frame}{Q333 --- Synthetic Data Trade-offs}
	\begin{block}{Question}
		Which of the following best describes a trade-off when using synthetic data?
		\begin{enumerate}
			\item It is always superior to real data.
			\item It reduces privacy risk and can be designed to remove bias, but may not capture the full complexity of real data.
			\item It eliminates the need for validation.
			\item It is always cheaper to generate.
	\end{enumerate}  \end{block}
	\begin{exampleblock}{Answer: B}\end{exampleblock}
	\begin{block}{Explanation}
		Synthetic data reduces privacy risk but may lack realism. Options A, C, and D are overstatements.
	\end{block}
\end{frame}

\begin{frame}{Q334 --- Fairness Trade-offs}
	\begin{block}{Question}
		Why must an organisation make a normative choice among fairness metrics?
		\begin{enumerate}
			\item Because all metrics agree.
			\item Because different metrics can conflict when base rates differ, so the choice reflects values and business context.
			\item Because fairness is a technical problem.
			\item Because regulators prescribe one metric.
	\end{enumerate}  \end{block}
	\begin{exampleblock}{Answer: B}\end{exampleblock}
	\begin{block}{Explanation}
		The impossibility result forces a choice among fairness criteria, reflecting normative considerations. Options A, C, and D are incorrect.
	\end{block}
\end{frame}

\begin{frame}{Q335 --- Local vs Global Explainability}
	\begin{block}{Question}
		Which is an example of local explainability?
		\begin{enumerate}
			\item Feature importance for the whole model.
			\item SHAP values for a specific prediction.
			\item Partial dependence plot.
			\item Global surrogate model.
	\end{enumerate}  \end{block}
	\begin{exampleblock}{Answer: B}\end{exampleblock}
	\begin{block}{Explanation}
		SHAP for a single prediction is local. Options A, C, and D are typically global.
	\end{block}
\end{frame}

\begin{frame}{Q336 --- Explainability in Practice}
	\begin{block}{Question}
		A company wants to give actionable advice to rejected applicants. Which technique is most appropriate?
		\begin{enumerate}
			\item Global feature importance.
			\item Counterfactual explanations (e.g., ``If income were \$X higher, you would be approved'').
			\item Model training accuracy.
			\item Model architecture description.
	\end{enumerate}  \end{block}
	\begin{exampleblock}{Answer: B}\end{exampleblock}
	\begin{block}{Explanation}
		Counterfactuals provide actionable insights. Options A, C, and D are not actionable for the individual.
	\end{block}
\end{frame}

\begin{frame}{Q337 --- Fairness Automation Limits}
	\begin{block}{Question}
		Why can't fairness be fully automated?
		\begin{enumerate}
			\item It is too expensive.
			\item Because it involves normative choices among conflicting criteria, especially when base rates differ.
			\item Because there is not enough data.
			\item Because GDPR forbids it.
	\end{enumerate}  \end{block}
	\begin{exampleblock}{Answer: B}\end{exampleblock}
	\begin{block}{Explanation}
		Fairness involves value judgments. Options A, C, and D are not the fundamental reason.
	\end{block}
\end{frame}

\begin{frame}{Q338 --- Global Inequality}
	\begin{block}{Question}
		Which of the following is a legitimate concern about AI and global inequality?
		\begin{enumerate}
			\item AI is naturally equitable.
			\item The concentration of AI capability in a few countries could worsen global disparities.
			\item AI is too slow to matter.
			\item AI reduces inequality automatically.
	\end{enumerate}  \end{block}
	\begin{exampleblock}{Answer: B}\end{exampleblock}
	\begin{block}{Explanation}
		Concentration of AI capability can widen global gaps. Options A, C, and D are incorrect or unsupported.
	\end{block}
\end{frame}

\begin{frame}{Q339 --- Privacy vs Utility}
	\begin{block}{Question}
		Which best describes the privacy-utility trade-off in AI?
		\begin{enumerate}
			\item More data always improves privacy.
			\item More data can improve model utility but increases privacy risk; techniques like differential privacy can help balance.
			\item Privacy and utility never conflict.
			\item Privacy is always more important.
	\end{enumerate}  \end{block}
	\begin{exampleblock}{Answer: B}\end{exampleblock}
	\begin{block}{Explanation}
		Data utility comes at privacy cost; techniques like differential privacy manage the trade-off. Options A, C, and D are oversimplifications.
	\end{block}
\end{frame}

\begin{frame}{Q340 --- AI Ethics Frameworks Benefit}
	\begin{block}{Question}
		The primary benefit of an AI ethics framework is:
		\begin{enumerate}
			\item Increasing model accuracy.
			\item Providing structured guidance for ethical decisions and reducing regulatory and reputational risk.
			\item Reducing training time.
			\item Eliminating all ethical issues.
	\end{enumerate}  \end{block}
	\begin{exampleblock}{Answer: B}\end{exampleblock}
	\begin{block}{Explanation}
		Ethics frameworks provide guidance and risk reduction. Options A, C, and D are not their primary purpose.
	\end{block}
\end{frame}

% ---------------- MODULE 5 — DATA AND AI MODEL GOVERNANCE ----------------

\begin{frame}{Q341 --- Data Governance Framework}
	\begin{block}{Question}
		Which of the following is NOT typically a component of a data governance framework?
		\begin{enumerate}
			\item Data strategy and quality.
			\item Data provenance and classification.
			\item Metadata management.
			\item Marketing strategy.
	\end{enumerate}  \end{block}
	\begin{exampleblock}{Answer: D}\end{exampleblock}
	\begin{block}{Explanation}
		Data governance covers strategy, quality, provenance, classification, metadata, protection, access, and compliance. Marketing is separate.
	\end{block}
\end{frame}

\begin{frame}{Q342 --- Data Provenance}
	\begin{block}{Question}
		Data provenance refers to:
		\begin{enumerate}
			\item The data storage location.
			\item The origin, history, and legal right to use data.
			\item The data transformation steps only.
			\item The data format.
	\end{enumerate}  \end{block}
	\begin{exampleblock}{Answer: B}\end{exampleblock}
	\begin{block}{Explanation}
		Provenance tracks origin, ownership, and legal rights. Options A, C, and D are partial or incorrect.
	\end{block}
\end{frame}

\begin{frame}{Q343 --- Data Classification}
	\begin{block}{Question}
		Which of the following best describes data classification?
		\begin{enumerate}
			\item Sorting data alphabetically.
			\item Categorising data by sensitivity and structure to determine handling and controls.
			\item Storing data in a database.
			\item Compressing data.
	\end{enumerate}  \end{block}
	\begin{exampleblock}{Answer: B}\end{exampleblock}
	\begin{block}{Explanation}
		Classification determines appropriate handling and protections. Options A, C, and D are unrelated.
	\end{block}
\end{frame}

\begin{frame}{Q344 --- Metadata Management}
	\begin{block}{Question}
		Metadata management is important because:
		\begin{enumerate}
			\item Metadata replaces data.
			\item Metadata supports discovery, lineage, governance, and understanding of data assets.
			\item Metadata is encrypted data.
			\item Metadata is unnecessary.
	\end{enumerate}  \end{block}
	\begin{exampleblock}{Answer: B}\end{exampleblock}
	\begin{block}{Explanation}
		Metadata enables discovery, lineage, and governance. Options A, C, and D misstate or dismiss.
	\end{block}
\end{frame}

\begin{frame}{Q345 --- Data Access Controls}
	\begin{block}{Question}
		Role-Based Access Control (RBAC) grants access based on:
		\begin{enumerate}
			\item Individual discretion.
			\item User roles within the organisation.
			\item Random selection.
			\item Data size.
	\end{enumerate}  \end{block}
	\begin{exampleblock}{Answer: B}\end{exampleblock}
	\begin{block}{Explanation}
		RBAC grants access based on roles. Options A, C, and D are not RBAC.
	\end{block}
\end{frame}

\begin{frame}{Q346 --- Model Governance Framework}
	\begin{block}{Question}
		A comprehensive model governance framework should include:
		\begin{enumerate}
			\item Definitions of models and non-models, risk tiering, validation standards, documentation, monitoring, and reporting.
			\item Only model training.
			\item Only model deployment.
			\item Only model retirement.
	\end{enumerate}  \end{block}
	\begin{exampleblock}{Answer: A}\end{exampleblock}
	\begin{block}{Explanation}
		Governance covers definitions, tiering, validation, documentation, monitoring, and reporting. Options B, C, and D are partial.
	\end{block}
\end{frame}

\begin{frame}{Q347 --- Model Documentation}
	\begin{block}{Question}
		The primary purpose of model documentation is to:
		\begin{enumerate}
			\item Improve accuracy.
			\item Provide a record of design, assumptions, limitations, and usage to support validation, audit, and governance.
			\item Speed training.
			\item Reduce features.
	\end{enumerate}  \end{block}
	\begin{exampleblock}{Answer: B}\end{exampleblock}
	\begin{block}{Explanation}
		Documentation supports validation, audit, and governance. Options A, C, and D are incorrect purposes.
	\end{block}
\end{frame}

\begin{frame}{Q348 --- Model Inventory}
	\begin{block}{Question}
		A model inventory should typically include:
		\begin{enumerate}
			\item Model owner, purpose, risk tier, validation status, and interdependencies.
			\item Only model names.
			\item Only model versions.
			\item Only model code.
	\end{enumerate}  \end{block}
	\begin{exampleblock}{Answer: A}\end{exampleblock}
	\begin{block}{Explanation}
		A model inventory is comprehensive: owner, purpose, risk tier, validation, and interdependencies. Options B, C, and D are incomplete.
	\end{block}
\end{frame}

\begin{frame}{Q349 --- Model Landscape}
	\begin{block}{Question}
		A model landscape differs from a model inventory in that it:
		\begin{enumerate}
			\item Lists all models.
			\item Shows interactions and interdependencies between models, highlighting cascade risks.
			\item Lists model code.
			\item Replaces the inventory.
	\end{enumerate}  \end{block}
	\begin{exampleblock}{Answer: B}\end{exampleblock}
	\begin{block}{Explanation}
		Landscape shows interactions; inventory is a registry. Options A, C, and D are incorrect.
	\end{block}
\end{frame}

\begin{frame}{Q350 --- Model Risk Management Roles}
	\begin{block}{Question}
		In the Three Lines of Defence model for AI/ML, who is primarily responsible for independent validation?
		\begin{enumerate}
			\item First line (model developers).
			\item Second line (independent risk management/model validation).
			\item Third line (internal audit).
			\item External auditors.
	\end{enumerate}  \end{block}
	\begin{exampleblock}{Answer: B}\end{exampleblock}
	\begin{block}{Explanation}
		The second line performs independent validation. First line owns; third line audits; external auditors add extra assurance. Options A, C, and D misstate the roles.
	\end{block}
\end{frame}

\end{document}